%---------------------------------------------------------------
\chapter{Penutup}
%---------------------------------------------------------------

%---------------------------------------------------------------
\section{Kesimpulan}
%---------------------------------------------------------------

Berdasarkan hasil penelitian yang telah dilakukan, dapat ditarik kesimpulan sebagai berikut.

\begin{enumerate}
    \item Model \textit{Joint Document-Level Sentiment-Topic Model} (JDST-DEC) berhasil dikembangkan dengan mengintegrasikan
     analisis sentimen dan pemodelan topik secara \textit{end-to-end} dalam satu kerangka \textit{multi-task learning}.
      Model memanfaatkan satu \textit{shared encoder} yang dipelajari secara bersama melalui 
      optimisasi \textit{sentiment loss} dan \textit{clustering loss}, sehingga mampu menghasilkan representasi 
      laten yang digunakan secara simultan oleh \textit{sentiment head} dan \textit{topic head} berbasis \textit{Deep Embedded Clustering} (DEC).

    \item Penerapan \textit{Deep Embedded Clustering} (DEC) pada \textit{topic head} berhasil membentuk representasi 
    topik dari representasi laten yang dihasilkan oleh \textit{shared encoder}. DEC mereduksi representasi BERT dari 
    768 dimensi menjadi 256 dimensi menggunakan \textit{autoencoder}, kemudian melakukan pembentukan klaster melalui 
    mekanisme \textit{soft assignment}. Pada evaluasi kualitas topik, JDST-DEC memperoleh nilai TC-NPMI tertinggi 
    sebesar 0{,}0803 pada dataset IMDB dan 0{,}0788 pada dataset Yelp. Hasil tersebut menunjukkan bahwa DEC mampu
     diintegrasikan dalam kerangka \textit{joint modelling}, meskipun kualitas topik yang dihasilkan masih dapat ditingkatkan.

    \item Pada tugas analisis sentimen, JDST-DEC memberikan performa yang lebih baik dibandingkan metode sebelumnya.
     Model secara konsisten memperoleh nilai \textit{accuracy} dan F1-score yang lebih tinggi dibandingkan NN-BERT 
     pada seluruh konfigurasi jumlah topik ($K=30$, $50$, dan $80$). Nilai F1-score terbaik yang diperoleh 
     mencapai 0{,}8986 pada dataset IMDB dan 0{,}9020 pada dataset Yelp, lebih tinggi dibandingkan NN-BERT yang 
     masing-masing memperoleh 0{,}8612 dan 0{,}8715. Hasil tersebut menunjukkan bahwa penerapan \textit{multi-task learning} 
     tidak menurunkan performa analisis sentimen, tetapi justru memberikan peningkatan dibandingkan model \textit{single-task}.

    \item Pada tugas pemodelan topik, JDST-DEC menunjukkan kualitas topik yang lebih baik dibandingkan MTL-GUI pada
     beberapa konfigurasi jumlah topik, namun masih berada di bawah BERTopic-KMeans berdasarkan metrik TC-NPMI maupun \textit{Topic Diversity}. 
     BERTopic-KMeans memperoleh TC-NPMI tertinggi sebesar 0{,}1874 pada IMDB dan 0{,}1465 pada Yelp, sedangkan nilai \textit{Topic Diversity} 
     tertinggi mencapai 0{,}9567 pada IMDB dan 0{,}7800 pada Yelp. Hasil tersebut menunjukkan bahwa meskipun DEC berhasil diintegrasikan 
     ke dalam model JDST, kualitas representasi topik yang dihasilkan masih belum mampu melampaui metode pemodelan topik khusus.
\end{enumerate}


%---------------------------------------------------------------
\section{Saran}
%---------------------------------------------------------------

Berdasarkan hasil dan keterbatasan pada penelitian ini, beberapa saran untuk
penelitian selanjutnya adalah sebagai berikut.

\begin{enumerate}
   \item Mengeksplorasi pengembangan arsitektur \textit{topic head} berbasis
    \textit{Deep Embedded Clustering} untuk meningkatkan kualitas representasi
    topik. Pengembangan dapat dilakukan melalui modifikasi struktur
    \textit{autoencoder}, mekanisme pembentukan distribusi target, maupun strategi
    optimisasi klaster.

    \item Mengeksplorasi pengembangan arsitektur \textit{sentiment head} untuk
    mendukung keseimbangan proses optimisasi antara tugas analisis sentimen
    dan klasterisasi. Berdasarkan hasil analisis \textit{multiloss}, perubahan
    pada komponen ini berpotensi membantu model mencapai keseimbangan yang
    lebih baik antara kualitas prediksi sentimen dan kualitas klaster.

    \item Mengeksplorasi strategi pembobotan \textit{loss} yang adaptif antara
    $\mathcal{L}_{\text{sent}}$ dan $\mathcal{L}_{\text{clust}}$ guna
    mengurangi \textit{trade-off} antara performa analisis sentimen dan
    kualitas klaster yang dihasilkan.
\end{enumerate}


